\documentclass[11pt]{article}
\usepackage[T1]{fontenc}
\usepackage[utf8]{inputenc}
\usepackage[ngerman]{babel}
\usepackage{amsmath,amssymb}
\usepackage[a4paper,margin=1in]{geometry}
\usepackage{microtype}
\usepackage{enumitem}
\setlist{nosep}
\newcommand{\RR}{\mathfrak{R}}
\begin{document}

\begin{center}
{\Large L.}\\[0.4em]
{\LARGE Stetigkeit und irrationale Zahlen}\\[0.6em]
{[Erste Auflage 1872, Fünfte Auflage 1927.]}\\[2em]
{\itshape Seinem geliebten Vater,}\\[0.4em]
{dem}\\[0.4em]
{\itshape Geh. Hofrat, Professor, Dr. jur. Julius Levin Ulrich Dedekind}\\[0.4em]
{\itshape in Braunschweig bei Gelegenheit seines fünfzigjährigen Amts-Jubiläums}\\[0.4em]
{\itshape am 26. April 1872 gewidmet.}
\end{center}

\bigskip

\noindent\textbf{Vorwort.}\dotfill 315\\
\textbf{\S 1. Eigenschaften der rationalen Zahlen.}\dotfill 317\\
\textbf{\S 2. Vergleichung der rationalen Zahlen mit den Punkten einer geraden Linie.}\dotfill 319\\
\textbf{\S 3. Stetigkeit der geraden Linie.}\dotfill 320\\
\textbf{\S 4. Schöpfung der irrationalen Zahlen.}\dotfill 323

\section*{Vorwort}

Die Betrachtungen, welche den Gegenstand dieser kleinen Schrift bilden, stammen aus dem Herbst des Jahres 1858. Ich befand mich damals als Professor am eidgenössischen Polytechnikum zu Zürich zum ersten Male in der Lage, die Elemente der Differentialrechnung vortragen zu müssen, und fühlte dabei empfindlicher als jemals früher den Mangel einer wirklich wissenschaftlichen Begründung der Arithmetik. Bei dem Begriffe der Annäherung einer veränderlichen Größe an einen festen Grenzwert und namentlich bei dem Beweise des Satzes, daß jede Größe, welche beständig, aber nicht über alle Grenzen wächst, sich gewiß einem Grenzwert nähern muß, nahm ich meine Zuflucht zu geometrischen Evidenzen. Auch jetzt halte ich ein solches Heranziehen geometrischer Anschauung bei dem ersten Unterrichte in der Differentialrechnung vom didaktischen Standpunkte aus für außerordentlich nützlich, ja unentbehrlich, wenn man nicht gar zu viel Zeit verlieren will. Aber daß diese Art der Einführung in die Differentialrechnung keinen Anspruch auf Wissenschaftlichkeit machen kann, wird wohl niemand leugnen. Für mich war damals dies Gefühl der Unbefriedigung ein so überwältigendes, daß ich den festen Entschluß faßte, so lange nachzudenken, bis ich eine rein arithmetische und völlig strenge Begründung der Prinzipien der Infinitesimalanalysis gefunden haben würde. Man sagt so häufig, die Differentialrechnung beschäftige sich mit den stetigen Größen, und doch wird nirgends eine Erklärung von dieser Stetigkeit gegeben, und auch die strengsten Darstellungen der Differentialrechnung gründen ihre Beweise nicht auf die Stetigkeit, sondern sie appellieren entweder mit mehr oder weniger Bewußtsein an geometrische, oder durch die Geometrie veranlaßte Vorstellungen, oder aber sie stützen sich auf solche Sätze, welche selbst nie rein arithmetisch bewiesen sind. Zu diesen gehört z. B. der oben erwähnte Satz, und eine genauere Untersuchung überzeugte mich, daß dieser oder auch jeder mit ihm äquivalente Satz gewissermaßen als ein hinreichendes Fundament für die Infinitesimalanalysis angesehen werden kann. Es kam nur noch darauf an, seinen eigentlichen Ursprung in den Elementen der Arithmetik zu entdecken und hiermit zugleich eine wirkliche Definition von dem Wesen der Stetigkeit zu gewinnen. Dies gelang mir am 24. November 1858, und wenige Tage darauf teilte ich das Ergebnis meines Nachdenkens meinem teuren Freunde Durege mit, was zu einer langen und lebhaften Unterhaltung führte. Später habe ich wohl dem einen oder anderen meiner Schüler diese Gedanken über eine wissenschaftliche Begründung der Arithmetik auseinandergesetzt, auch hier in Braunschweig in dem wissenschaftlichen Verein der Professoren einen Vortrag über diesen Gegenstand gehalten, aber zu einer eigentlichen Publikation konnte ich mich nicht recht entschließen, weil erstens die Darstellung nicht ganz leicht, und weil außerdem die Sache so wenig fruchtbar ist. Indessen hatte ich doch schon halb und halb daran gedacht, dieses Thema zum Gegenstände dieser Gelegenheitsschrift zu wählen, als vor wenigen Tagen, am 14. März, die Abhandlung: „Die Elemente der Funktionenlehre“, von E. Heine (Crelle's Journal, Bd. 74) durch die Güte ihres hochverehrten Verfassers in meine Hände gelangte und mich in meinem Entschlüsse bestärkte. Dem Wesen nach stimme ich zwar vollständig mit dem Inhalte dieser Schrift überein, wie es ja nicht anders sein kann, aber ich will freimütig gestehen, daß meine Darstellung mir der Form nach einfacher zu sein und den eigentlichen Kernpunkt präziser hervorzuheben scheint. Und während ich an diesem Vorwort schreibe (20. März 1872), erhalte ich die interessante Abhandlung: „Über die Ausdehnung eines Satzes aus der Theorie der trigonometrischen Reihen“, von G. Cantor (Math. Annalen von Clebsch und Neumann, Bd. 5), für welche ich dem scharfsinnigen Verfasser meinen besten Dank sage. Wie ich bei raschem Durchlesen finde, so stimmt das Axiom in \S 2 derselben, abgesehen von der äußeren Form der Einkleidung, vollständig mit dem überein, was ich unten in \S 3 als das Wesen der Stetigkeit bezeichne. Welchen Nutzen aber die wenn auch nur begriffliche Unterscheidung von reellen Zahlgrößen noch höherer Art gewähren wird, vermag ich gerade nach meiner Auffassung des in sich vollkommenen reellen Zahlgebietes noch nicht zu erkennen.

\section*{\S 1. Eigenschaften der rationalen Zahlen}

Die Entwicklung der Arithmetik der rationalen Zahlen wird hier zwar vorausgesetzt, doch halte ich es für gut, einige Hauptmomente ohne Diskussion hervorzuheben, nur um den Standpunkt von vornherein zu bezeichnen, den ich im folgenden einnehme. Ich sehe die ganze Arithmetik als eine notwendige oder wenigstens natürliche Folge des einfachsten arithmetischen Aktes, des Zählens, an, und das Zählen selbst ist nichts anderes als die sukzessive Schöpfung der unendlichen Reihe der positiven ganzen Zahlen, in welcher jedes Individuum durch das unmittelbar vorhergehende definiert wird; der einfachste Akt ist der Übergang von einem schon erschaffenen Individuum zu dem darauffolgenden neu zu erschaffenden. Die Kette dieser Zahlen bildet an sich schon ein überaus nützliches Hilfsmittel für den menschlichen Geist, und sie bietet einen unerschöpflichen Reichtum an merkwürdigen Gesetzen dar, zu welchen man durch die Einführung der vier arithmetischen Grundoperationen gelangt. Die Addition ist die Zusammenfassung einer beliebigen Wiederholung des obigen einfachsten Aktes zu einem einzigen Akte, und aus ihr entspringt auf dieselbe Weise die Multiplikation. Während diese beiden Operationen stets ausführbar sind, zeigen die umgekehrten Operationen, die Subtraktion und Division, nur eine beschränkte Zulässigkeit. Welches nun auch die nächste Veranlassung gewesen sein mag, welche Vergleichungen oder Analogieen mit Erfahrungen, Anschauungen dazu geführt haben mögen, bleibe dahingestellt; genug, gerade diese Beschränktheit in der Ausführbarkeit der indirekten Operationen ist jedesmal die eigentliche Ursache eines neuen Schöpfungsaktes geworden; so sind die negativen und gebrochenen Zahlen durch den menschlichen Geist erschaffen, und es ist in dem System aller rationalen Zahlen ein Instrument von unendlich viel größerer Vollkommenheit gewonnen. Dieses System, welches ich mit $R$ bezeichnen will, besitzt vor allen Dingen eine Vollständigkeit und Abgeschlossenheit, welche ich an einem anderen Orte\footnote{Vorlesungen über Zahlentheorie von P. G. Lejeune Dirichlet. Zweite Auflage. \S 159.} als Merkmal eines Zahlkörpers bezeichnet habe, und welche darin besteht, daß die vier Grundoperationen mit je zwei Individuen in $R$ stets ausführbar sind, d. h. daß das Resultat derselben stets wieder ein bestimmtes Individuum in $R$ ist, wenn man den einzigen Fall der Division durch die Zahl Null ausnimmt.

Für unseren nächsten Zweck ist aber noch wichtiger eine andere Eigenschaft des Systems $R$, welche man dahin aussprechen kann, daß das System $R$ ein wohlgeordnetes, nach zwei entgegengesetzten Seiten hin unendliches Gebiet von einer Dimension bildet. Was damit gemeint sein soll, ist durch die Wahl der Ausdrücke, welche geometrischen Vorstellungen entlehnt sind, hinreichend angedeutet; um so notwendiger ist es, die entsprechenden rein arithmetischen Eigentümlichkeiten hervorzuheben, damit es auch nicht einmal den Anschein behält, als bedürfte die Arithmetik solcher ihr fremden Vorstellungen.

Soll ausgedrückt werden, daß die Zeichen $a$ und $b$ eine und dieselbe rationale Zahl bedeuten, so setzt man sowohl $a=b$ wie $b=a$. Die Verschiedenheit zweier rationaler Zahlen $a,b$ zeigt sich darin, daß die Differenz $a-b$ entweder einen positiven oder einen negativen Wert hat. Im ersten Falle heißt $a$ größer als $b$, $b$ kleiner als $a$, was auch durch die Zeichen $a>b$, $b<a$ angedeutet wird.\footnote{Es ist also im folgenden immer das sogenannte „algebraische“ größer und kleiner Sein gemeint, wenn nicht das Wort „absolut“ hinzugefügt wird.} Da im zweiten Falle $b-a$ einen positiven Wert hat, so ist $b>a$, $a<b$. Hinsichtlich dieser doppelten Möglichkeit in der Art der Verschiedenheit gelten nun folgende Gesetze.

\begin{enumerate}[label=\Roman*.]
\item Ist $a>b$, und $b>c$, so ist $a>c$. Wir wollen jedesmal, wenn $a,c$ zwei verschiedene (oder ungleiche) Zahlen sind, und wenn $b$ größer als die eine, kleiner als die andere ist, ohne Scheu vor dem Anklang an geometrische Vorstellungen dies kurz so ausdrücken: $b$ liegt zwischen den beiden Zahlen $a,c$.
\item Sind $a,c$ zwei verschiedene Zahlen, so gibt es immer unendlich viele verschiedene Zahlen $b$, welche zwischen $a,c$ liegen.
\item Ist $a$ eine bestimmte Zahl, so zerfallen alle Zahlen des Systems $R$ in zwei Klassen $A_1$ und $A_2$, deren jede unendlich viele Individuen enthält; die erste Klasse $A_1$ umfaßt alle Zahlen $a_1$, welche $<a$ sind, die zweite Klasse $A_2$ umfaßt alle Zahlen $a_2$, welche $>a$ sind; die Zahl $a$ selbst kann nach Belieben der ersten oder der zweiten Klasse zugeteilt werden, und sie ist dann entsprechend die größte Zahl der ersten oder die kleinste Zahl der zweiten Klasse. In jedem Falle ist die Zerlegung des Systems $R$ in die beiden Klassen $A_1,A_2$ von der Art, daß jede Zahl der ersten Klasse $A_1$ kleiner als jede Zahl der zweiten Klasse $A_2$ ist.
\end{enumerate}

\section*{\S 2. Vergleichung der rationalen Zahlen mit den Punkten einer geraden Linie}

Die soeben hervorgehobenen Eigenschaften der rationalen Zahlen erinnern an die gegenseitigen Lagenbeziehungen zwischen den Punkten einer geraden Linie $L$. Werden die beiden in ihr existierenden entgegengesetzten Richtungen durch „rechts“ und „links“ unterschieden, und sind $p,q$ zwei verschiedene Punkte, so liegt entweder $p$ rechts von $q$, und gleichzeitig $q$ links von $p$, oder umgekehrt, es liegt $q$ rechts von $p$, und gleichzeitig $p$ links von $q$. Ein dritter Fall ist unmöglich, wenn $p,q$ wirklich verschiedene Punkte sind. Hinsichtlich dieser Lagenverschiedenheit bestehen folgende Gesetze.

\begin{enumerate}[label=\Roman*.]
\item Liegt $p$ rechts von $q$, und $q$ wieder rechts von $r$, so liegt auch $p$ rechts von $r$; und man sagt, daß $q$ zwischen den Punkten $p$ und $r$ liegt.
\item Sind $p,r$ zwei verschiedene Punkte, so gibt es immer unendlich viele Punkte $q$, welche zwischen $p$ und $r$ liegen.
\item Ist $p$ ein bestimmter Punkt in $L$, so zerfallen alle Punkte in $L$ in zwei Klassen, $P_1,P_2$, deren jede unendlich viele Individuen enthält; die erste Klasse $P_1$ umfaßt alle die Punkte $p_1$, welche links von $p$ liegen, und die zweite Klasse $P_2$ umfaßt alle die Punkte $p_2$, welche rechts von $p$ liegen; der Punkt $p$ selbst kann nach Belieben der ersten oder der zweiten Klasse zugeteilt werden. In jedem Falle ist die Zerlegung der Geraden $L$ in die beiden Klassen oder Stücke $P_1,P_2$ von der Art, daß jeder Punkt der ersten Klasse $P_1$ links von jedem Punkte der zweiten Klasse $P_2$ liegt.
\end{enumerate}

Diese Analogie zwischen den rationalen Zahlen und den Punkten einer Geraden wird bekanntlich zu einem wirklichen Zusammenhange, wenn in der Geraden ein bestimmter Anfangspunkt oder Nullpunkt $o$ und eine bestimmte Längeneinheit zur Ausmessung der Strecken gewählt wird. Mit Hilfe der letzteren kann für jede rationale Zahl $a$ eine entsprechende Länge konstruiert werden, und trägt man dieselbe von dem Punkte $o$ aus nach rechts oder links auf der Geraden ab, je nachdem $a$ positiv oder negativ ist, so gewinnt man einen bestimmten Endpunkt $p$, welcher als der der Zahl $a$ entsprechende Punkt bezeichnet werden kann; der rationalen Zahl Null entspricht der Punkt $o$. Auf diese Weise entspricht jeder rationalen Zahl $a$, d. h. jedem Individuum in $R$, ein und nur ein Punkt $p$, d. h. ein Individuum in $L$. Entsprechen den beiden Zahlen $a,b$ bzw. die beiden Punkte $p,q$, und ist $a>b$, so liegt $p$ rechts von $q$. Den Gesetzen I, II, III des vorigen Paragraphen entsprechen vollständig die Gesetze I, II, III des jetzigen.

\section*{\S 3. Stetigkeit der geraden Linie}

Von der größten Wichtigkeit ist nun aber die Tatsache, daß es in der Geraden $L$ unendlich viele Punkte gibt, welche keiner rationalen Zahl entsprechen. Entspricht nämlich der Punkt $p$ der rationalen Zahl $a$, so ist bekanntlich die Länge $op$ kommensurabel mit der bei der Konstruktion benutzten unabänderlichen Längeneinheit, d. h. es gibt eine dritte Länge, ein sogenanntes gemeinschaftliches Maß, von welcher diese beiden Längen ganze Vielfache sind. Aber schon die alten Griechen haben gewußt und bewiesen, daß es Längen gibt, welche mit einer gegebenen Längeneinheit inkommensurabel sind, z. B. die Diagonale des Quadrates, dessen Seite die Längeneinheit ist. Trägt man eine solche Länge von dem Punkt $o$ aus auf der Geraden ab, so erhält man einen Endpunkt, welcher keiner rationalen Zahl entspricht. Da sich ferner leicht beweisen läßt, daß es unendlich viele Längen gibt, welche mit der Längeneinheit inkommensurabel sind, so können wir behaupten: Die Gerade $L$ ist unendlich viel reicher an Punktindividuen, als das Gebiet $R$ der rationalen Zahlen an Zahlindividuen.

Will man nun, was doch der Wunsch ist, alle Erscheinungen in der Geraden auch arithmetisch verfolgen, so reichen dazu die rationalen Zahlen nicht aus, und es wird daher unumgänglich notwendig, das Instrument $R$, welches durch die Schöpfung der rationalen Zahlen konstruiert war, wesentlich zu verfeinern durch eine Schöpfung von neuen Zahlen der Art, daß das Gebiet der Zahlen dieselbe Vollständigkeit oder, wie wir gleich sagen wollen, dieselbe Stetigkeit gewinnt, wie die gerade Linie.

Die bisherigen Betrachtungen sind allen so bekannt und geläufig, daß viele ihre Wiederholung für sehr überflüssig erachten werden. Dennoch hielt ich diese Rekapitulation für notwendig, um die Hauptfrage gehörig vorzubereiten. Die bisher übliche Einführung der irrationalen Zahlen knüpft nämlich geradezu an den Begriff der extensiven Größen an — welcher aber selbst nirgends streng definiert wird — und erklärt die Zahl als das Resultat der Messung einer solchen Größe durch eine zweite gleichartige.\footnote{Der scheinbare Vorzug der Allgemeinheit dieser Definition der Zahl schwindet sofort dahin, wenn man an die komplexen Zahlen denkt. Nach meiner Auffassung kann umgekehrt der Begriff des Verhältnisses zwischen zwei gleichartigen Größen erst dann klar entwickelt werden, wenn die irrationalen Zahlen schon eingeführt sind.} Statt dessen fordere ich, daß die Arithmetik sich aus sich selbst heraus entwickeln soll. Daß solche Anknüpfungen an nicht arithmetische Vorstellungen die nächste Veranlassung zur Erweiterung des Zahlbegriffes gegeben haben, mag im allgemeinen zugegeben werden (doch ist dies bei der Einführung der komplexen Zahlen entschieden nicht der Fall gewesen); aber hierin liegt ganz gewiß kein Grund, diese fremdartigen Betrachtungen selbst in die Arithmetik, in die Wissenschaft von den Zahlen aufzunehmen. So wie die negativen und gebrochenen rationalen Zahlen durch eine freie Schöpfung hergestellt, und wie die Gesetze der Rechnungen mit diesen Zahlen auf die Gesetze der Rechnungen mit ganzen positiven Zahlen zurückgeführt werden müssen und können, ebenso hat man dahin zu streben, daß auch die irrationalen Zahlen durch die rationalen Zahlen allein vollständig definiert werden. Nur das Wie? bleibt die Frage.

Die obige Vergleichung des Gebietes $R$ der rationalen Zahlen mit einer Geraden hat zu der Erkenntnis der Lückenhaftigkeit, Unvollständigkeit oder Unstetigkeit des ersteren geführt, während wir der Geraden Vollständigkeit, Lückenlosigkeit oder Stetigkeit zuschreiben. Worin besteht denn nun eigentlich diese Stetigkeit? In der Beantwortung dieser Frage muß alles enthalten sein, und nur durch sie wird man eine wissenschaftliche Grundlage für die Untersuchung aller stetigen Gebiete gewinnen. Mit vagen Reden über den ununterbrochenen Zusammenhang in den kleinsten Teilen ist natürlich nichts erreicht; es kommt darauf an, ein präzises Merkmal der Stetigkeit anzugeben, welches als Basis für wirkliche Deduktionen gebraucht werden kann. Lange Zeit habe ich vergeblich darüber nachgedacht, aber endlich fand ich, was ich suchte. Dieser Fund wird von verschiedenen Personen vielleicht verschieden beurteilt werden, doch glaube ich, daß die meisten seinen Inhalt sehr trivial finden werden. Er besteht im folgenden. Im vorigen Paragraphen ist darauf aufmerksam gemacht, daß jeder Punkt $p$ der Geraden eine Zerlegung derselben in zwei Stücke von der Art hervorbringt, daß jeder Punkt des einen Stückes links von jedem Punkte des anderen liegt. Ich finde nun das Wesen der Stetigkeit in der Umkehrung, also in dem folgenden Prinzip:

\begin{quote}
„Zerfallen alle Punkte der Geraden in zwei Klassen von der Art, daß jeder Punkt der ersten Klasse links von jedem Punkte der zweiten Klasse liegt, so existiert ein und nur ein Punkt, welcher diese Einteilung aller Punkte in zwei Klassen, diese Zerschneidung der Geraden in zwei Stücke hervorbringt.“
\end{quote}

Wie schon gesagt, glaube ich nicht zu irren, wenn ich annehme, daß jedermann die Wahrheit dieser Behauptung sofort zugeben wird; die meisten meiner Leser werden sehr enttäuscht sein, zu vernehmen, daß durch diese Trivialität das Geheimnis der Stetigkeit enthüllt sein soll. Dazu bemerke ich folgendes. Es ist mir sehr lieb, wenn jedermann das obige Prinzip so einleuchtend findet und so übereinstimmend mit seinen Vorstellungen von einer Linie; denn ich bin außerstande, irgendeinen Beweis für seine Richtigkeit beizubringen, und niemand ist dazu imstande. Die Annahme dieser Eigenschaft der Linie ist nichts als ein Axiom, durch welches wir erst der Linie ihre Stetigkeit zuerkennen, durch welches wir die Stetigkeit in die Linie hineindenken. Hat überhaupt der Raum eine reale Existenz, so braucht er doch nicht notwendig stetig zu sein; unzählige seiner Eigenschaften würden dieselben bleiben, wenn er auch unstetig wäre. Und wüßten wir gewiß, daß der Raum unstetig wäre, so könnte uns doch wieder nichts hindern, falls es uns beliebte, ihn durch Ausfüllung seiner Lücken in Gedanken zu einem stetigen zu machen; diese Ausfüllung würde aber in einer Schöpfung von neuen Punktindividuen bestehen und dem obigen Prinzip gemäß auszuführen sein.

\section*{\S 4. Schöpfung der irrationalen Zahlen}

Durch die letzten Worte ist schon hinreichend angedeutet, auf welche Art das unstetige Gebiet $R$ der rationalen Zahlen zu einem stetigen vervollständigt werden muß. In \S 1 ist hervorgehoben (III), daß jede rationale Zahl $a$ eine Zerlegung des Systems $R$ in zwei Klassen $A_1,A_2$ von der Art hervorbringt, daß jede Zahl $a_1$ der ersten Klasse $A_1$ kleiner ist als jede Zahl $a_2$ der zweiten Klasse $A_2$; die Zahl $a$ ist entweder die größte Zahl der Klasse $A_1$, oder die kleinste Zahl der Klasse $A_2$. Ist nun irgendeine Einteilung des Systems $R$ in zwei Klassen $A_1,A_2$ gegeben, welche nur die charakteristische Eigenschaft besitzt, daß jede Zahl $a_1$ in $A_1$ kleiner ist als jede Zahl $a_2$ in $A_2$, so wollen wir der Kürze halber eine solche Einteilung einen Schnitt nennen und mit $(A_1,A_2)$ bezeichnen. Wir können dann sagen, daß jede rationale Zahl $a$ einen Schnitt oder eigentlich zwei Schnitte hervorbringt, welche wir aber nicht als wesentlich verschieden ansehen wollen; dieser Schnitt hat außerdem die Eigenschaft, daß entweder unter den Zahlen der ersten Klasse eine größte, oder unter den Zahlen der zweiten Klasse eine kleinste existiert. Und umgekehrt, besitzt ein Schnitt auch diese Eigenschaft, so wird er durch diese größte oder kleinste rationale Zahl hervorgebracht.

Aber man überzeugt sich leicht, daß auch unendlich viele Schnitte existieren, welche nicht durch rationale Zahlen hervorgebracht werden. Das nächstliegende Beispiel ist folgendes.

Es sei $D$ eine positive ganze Zahl, aber nicht das Quadrat einer ganzen Zahl, so gibt es eine positive ganze Zahl $\lambda$ von der Art, daß
\[
\lambda^2<D<(\lambda+1)^2
\]
wird.

Nimmt man in die zweite Klasse $A_2$ jede positive rationale Zahl $a_2$, deren Quadrat $>D$ ist, in die erste Klasse $A_1$ aber alle anderen rationalen Zahlen $a_1$, so bildet diese Einteilung einen Schnitt $(A_1,A_2)$, d. h. jede Zahl $a_1$ ist kleiner als jede Zahl $a_2$. Ist nämlich $a_1=0$ oder negativ, so ist $a_1$ schon aus diesem Grunde kleiner als jede Zahl $a_2$, weil diese zufolge der Definition positiv ist; ist aber $a_1$ positiv, so ist ihr Quadrat $<D$, und folglich ist $a_1$ kleiner als jede positive Zahl $a_2$, deren Quadrat $>D$ ist.

Dieser Schnitt wird aber durch keine rationale Zahl hervorgebracht. Um dies zu beweisen, muß vor allem gezeigt werden, daß es keine rationale Zahl gibt, deren Quadrat $=D$ ist. Obgleich dies aus den ersten Elementen der Zahlentheorie bekannt ist, so mag doch hier der folgende indirekte Beweis Platz finden. Gibt es eine rationale Zahl, deren Quadrat $=D$ ist, so gibt es auch zwei positive ganze Zahlen $t,u$, welche der Gleichung
\[
t^2-Du^2=0
\]
genügen, und man darf annehmen, daß $u$ die kleinste positive ganze Zahl ist, welche die Eigenschaft besitzt, daß ihr Quadrat durch Multiplikation mit $D$ in das Quadrat einer ganzen Zahl $t$ verwandelt wird. Da nun offenbar
\[
\lambda^2u^2<t^2<(\lambda+1)^2u^2
\]
ist, so wird die Zahl
\[
u'=t-\lambda u
\]
eine positive ganze Zahl, und zwar kleiner als $u$. Setzt man ferner
\[
t'=Du-\lambda t,
\]
so wird $t'$ ebenfalls eine positive ganze Zahl, und es ergibt sich
\[
t'^2-Du'^2=(\lambda^2-D)(t^2-Du^2)=0,
\]
was mit der Annahme über $u$ im Widerspruch steht.

Mithin ist das Quadrat einer jeden rationalen Zahl $x$ entweder $<D$ oder $>D$. Hieraus folgt nun leicht, daß es weder in der Klasse $A_1$ eine größte, noch in der Klasse $A_2$ eine kleinste Zahl gibt. Setzt man nämlich
\[
y=x\cdot \frac{x^2+3D}{3x^2+D},
\]
so ist
\[
y-x=\frac{2x(D-x^2)}{3x^2+D}
\]
und
\[
y^2-D=\frac{(x^2-D)^3}{(3x^2+D)^2}.
\]
Nimmt man hierin für $x$ eine positive Zahl aus der Klasse $A_1$, so ist $x^2<D$, und folglich wird $y>x$, und $y^2<D$, also gehört $y$ ebenfalls der Klasse $A_1$ an. Setzt man aber für $x$ eine Zahl aus der Klasse $A_2$, so ist $x^2>D$, und folglich wird $y<x$, $y$ positiv und $y^2>D$, also gehört $y$ ebenfalls der Klasse $A_2$ an. Dieser Schnitt wird daher durch keine rationale Zahl hervorgebracht.

In dieser Eigenschaft, daß nicht alle Schnitte durch rationale Zahlen hervorgebracht werden, besteht die Unvollständigkeit oder Unstetigkeit des Gebietes $R$ aller rationalen Zahlen.

Jedesmal nun, wenn ein Schnitt $(A_1,A_2)$ vorliegt, welcher durch keine rationale Zahl hervorgebracht wird, so erschaffen wir eine neue, eine irrationale Zahl $\alpha$, welche wir als durch diesen Schnitt $(A_1,A_2)$ vollständig definiert ansehen; wir werden sagen, daß die Zahl $\alpha$ diesem Schnitt entspricht, oder daß sie diesen Schnitt hervorbringt. Es entspricht also von jetzt ab jedem bestimmten Schnitt eine und nur eine bestimmte rationale oder irrationale Zahl und wir sehen zwei Zahlen stets und nur dann als verschieden oder ungleich an, wenn sie wesentlich verschiedenen Schnitten entsprechen.

Um nun eine Grundlage für die Anordnung aller reellen, d. h. aller rationalen und irrationalen Zahlen zu gewinnen, müssen wir zunächst die Beziehungen zwischen irgend zwei Schnitten $(A_1,A_2)$ und $(B_1,B_2)$ untersuchen, welche durch irgend zwei Zahlen $\alpha$ und $\beta$ hervorgebracht werden. Offenbar ist ein Schnitt $(A_1,A_2)$ schon vollständig gegeben, wenn eine der beiden Klassen, z. B. die erste $A_1$, bekannt ist, weil die zweite $A_2$ aus allen nicht in $A_1$ enthaltenen rationalen Zahlen besteht, und die charakteristische Eigenschaft einer solchen ersten Klasse $A_1$ liegt darin, daß sie, wenn die Zahl $a_1$ in ihr enthalten ist, auch alle kleineren Zahlen als $a_1$ enthält. Vergleicht man nun zwei solche erste Klassen $A_1,B_1$ miteinander, so kann es 1. sein, daß sie vollständig identisch sind, d. h. daß jede in $A_1$ enthaltene Zahl $a_1$ auch in $B_1$, und daß jede in $B_1$ enthaltene Zahl $b_1$ auch in $A_1$ enthalten ist. In diesem Falle ist dann notwendig auch $A_2$ identisch mit $B_2$, die beiden Schnitte sind vollständig identisch, was wir in Zeichen durch $\alpha=\beta$ oder $\beta=\alpha$ andeuten.

Sind aber die beiden Klassen $A_1,B_1$ nicht identisch, so gibt es in der einen, z. B. in $A_1$, eine Zahl $a_1'=b_2'$, welche nicht in der anderen $B_1$ enthalten ist, und welche sich folglich in $B_2$ vorfindet; mithin sind gewiß alle in $B_1$ enthaltenen Zahlen $b_1$ kleiner als diese Zahl $a_1'=b_2'$, und folglich sind alle Zahlen $b_1$ auch in $A_1$ enthalten.

Ist nun 2. diese Zahl $a_1'$ die einzige in $A_1$, welche nicht in $B_1$ enthalten ist, so ist jede andere in $A_1$ enthaltene Zahl $a_1$ in $B_1$ enthalten, und folglich kleiner als $a_1'$, d. h. $a_1'$ ist die größte unter allen Zahlen $a_1$, mithin wird der Schnitt $(A_1,A_2)$ durch die rationale Zahl $\alpha=a_1'=b_2'$ hervorgebracht. Von dem anderen Schnitte $(B_1,B_2)$ wissen wir schon, daß alle Zahlen $b_1$ in $B_1$ auch in $A_1$ enthalten und kleiner als die Zahl $a_1'=b_2'$ sind, welche in $B_2$ enthalten ist; jede andere in $B_2$ enthaltene Zahl $b_2$ muß aber größer als $a_1'=b_2'$ sein, weil sie sonst auch kleiner als $a_1'$, also in $A_1$ und folglich auch in $B_1$ enthalten wäre; mithin ist $b_2'$ die kleinste unter allen in $B_2$ enthaltenen Zahlen, und folglich wird auch der Schnitt $(B_1,B_2)$ durch dieselbe rationale Zahl $\beta=b_2'=a_1'=\alpha$ hervorgebracht. Die beiden Schnitte sind daher nur unwesentlich verschieden.

Gibt es aber 3. in $A_1$ wenigstens zwei verschiedene Zahlen $a_1'=b_2'$ und $a_1''=b_2''$, welche nicht in $B_1$ enthalten sind, so gibt es deren auch unendlich viele, weil alle die unendlich vielen zwischen $a_1'$ und $a_1''$ liegenden Zahlen (\S 1. II) offenbar in $A_1$, aber nicht in $B_1$ enthalten sind. In diesem Falle nennen wir die diesen beiden wesentlich verschiedenen Schnitten $(A_1,A_2)$ und $(B_1,B_2)$ entsprechenden Zahlen $\alpha$ und $\beta$ ebenfalls verschieden voneinander, und zwar sagen wir, daß $\alpha$ größer als $\beta$, daß $\beta$ kleiner als $\alpha$ ist, was wir in Zeichen sowohl durch $\alpha>\beta$, als durch $\beta<\alpha$ ausdrücken. Hierbei ist hervorzuheben, daß diese Definition vollständig mit der früheren zusammenfällt, wenn beide Zahlen $\alpha,\beta$ rational sind.

Die nun noch übrigen möglichen Fälle sind diese. Gibt es 4. in $B_1$ eine und nur eine Zahl $b_1'=a_2'$, welche nicht in $A_1$ enthalten ist, so sind die beiden Schnitte $(A_1,A_2)$ und $(B_1,B_2)$ nur unwesentlich verschieden und sie werden durch eine und dieselbe rationale Zahl $\alpha=a_2'=b_1'=\beta$ hervorgebracht. Gibt es aber 5. in $B_1$ mindestens zwei verschiedene Zahlen, welche nicht in $A_1$ enthalten sind, so ist $\beta>\alpha$, $\alpha<\beta$.

Da hiermit alle Fälle erschöpft sind, so ergibt sich, daß von zwei verschiedenen Zahlen notwendig die eine die größere, die andere die kleinere sein muß, was zwei Möglichkeiten enthält. Ein dritter Fall ist unmöglich. Dies lag zwar schon in der Wahl des Komparativs (größer, kleiner) zur Bezeichnung der Beziehung zwischen $\alpha,\beta$; aber diese Wahl ist erst jetzt nachträglich gerechtfertigt. Gerade bei solchen Untersuchungen hat man sich auf das sorgfältigste zu hüten, daß man selbst bei dem besten Willen, ehrlich zu sein, durch eine voreilige Wahl von Ausdrücken, welche anderen schon entwickelten Vorstellungen entlehnt sind, sich nicht verleiten lasse, unerlaubte Übertragungen aus dem einen Gebiete in das andere vorzunehmen.

Betrachtet man nun noch einmal genau den Fall $\alpha>\beta$, so ergibt sich, daß die kleinere Zahl $\beta$, wenn sie rational ist, gewiß der Klasse $A_1$ angehört; da es nämlich in $A_1$ eine Zahl $a_1'=b_2'$ gibt, welche der Klasse $B_2$ angehört, so ist die Zahl $\beta$, mag sie die größte Zahl in $B_1$ oder die kleinste Zahl in $B_2$ sein, gewiß $\leq a_1'$ und folglich in $A_1$ enthalten. Ebenso ergibt sich aus $\alpha>\beta$, daß die größere Zahl $\alpha$, wenn sie rational ist, gewiß der Klasse $B_2$ angehört, weil $\alpha\geq a_1'$ ist. Vereinigt man beide Betrachtungen, so erhält man folgendes Resultat: Wird ein Schnitt $(A_1,A_2)$ durch die Zahl $\alpha$ hervorgebracht, so gehört irgendeine rationale Zahl zu der Klasse $A_1$ oder zu der Klasse $A_2$, je nachdem sie kleiner oder größer ist als $\alpha$; ist die Zahl $\alpha$ selbst rational, so kann sie der einen oder der anderen Klasse angehören.

Hieraus ergibt sich endlich noch folgendes. Ist $\alpha>\beta$, gibt es also unendlich viele Zahlen in $A_1$, welche nicht in $B_1$ enthalten sind, so gibt es auch unendlich viele solche Zahlen, welche zugleich von $\alpha$ und von $\beta$ verschieden sind; jede solche rationale Zahl $c$ ist $<\alpha$, weil sie in $A_1$ enthalten ist, und sie ist zugleich $>\beta$, weil sie in $B_2$ enthalten ist.

\end{document}
